\documentclass[11pt]{article}

\usepackage[a4paper,margin=0.78in]{geometry}
\usepackage[T1]{fontenc}
\usepackage[utf8]{inputenc}
\usepackage{lmodern}
\usepackage{amsmath,amssymb,amsthm,mathtools}
\usepackage{enumitem,booktabs,array,tabularx,microtype,xcolor}
\usepackage[numbers,sort&compress]{natbib}
\usepackage[colorlinks=true,linkcolor=blue!65!black,
  citecolor=blue!65!black,urlcolor=blue!70!black]{hyperref}
\usepackage[nameinlink,noabbrev]{cleveref}
\setlength{\emergencystretch}{3em}

\hypersetup{
  pdftitle={MVP5 at the Frontier Occurrence Lift},
  pdfauthor={Liang Wang},
  pdfsubject={A source-locked route decision after the four-map and actual-core audits},
  pdfkeywords={MVP5, fixed shift, occurrence lift, route audit, arithmetic frontier}
}

\newtheorem{theorem}{Theorem}[section]
\newtheorem{proposition}[theorem]{Proposition}
\newtheorem{corollary}[theorem]{Corollary}
\theoremstyle{definition}
\newtheorem{definition}[theorem]{Definition}
\theoremstyle{remark}
\newtheorem{remark}[theorem]{Remark}

\newcommand{\Lzero}{\mathrm{L0}}
\newcommand{\Lone}{\mathrm{L1}}
\newcommand{\Ltwo}{\mathrm{L2}}
\newcommand{\GO}{\textnormal{\textsc{go}}}
\newcommand{\AINF}{\textnormal{\textsc{architecture-infeasible}}}
\newcommand{\REROUTE}{\textnormal{\textsc{reroute}}}
\newcommand{\STOPR}{\textnormal{\textsc{stop-route}}}
\newcommand{\NT}{\textnormal{\textsc{not-testable}}}
\newcommand{\AF}{\textnormal{\textsc{arithmetic-frontier}}}
\newcommand{\OPEN}{\textnormal{\textsc{open}}}
\newcommand{\INVALID}{\textnormal{\textsc{invalid}}}
\newcommand{\Bhard}{B_{h_0,\delta}}
\newcolumntype{Y}{>{\raggedright\arraybackslash}X}
\newcolumntype{P}[1]{>{\raggedright\arraybackslash}p{#1}}

\title{\textbf{MVP5 at the Frontier Occurrence Lift:}\\
\textbf{A Source-Locked Route Decision after the}\\
\textbf{Four-Map and Actual-Core Audits}}
\author{Liang Wang\\
School of Artificial Intelligence and Automation,\\
Huazhong University of Science and Technology, Wuhan 430074, P.R. China\\
\texttt{wangliang.f@gmail.com}}
\date{July 2026}

\begin{document}
\maketitle

\begin{abstract}
MVP4 located the first unsupported carrier after a complete cut.
TPC-143--151 now determine what that statement means path by path.
The four downstream maps share one first prerequisite: a conservative
row-separated occurrence lift from every eligible-tail and frontier
cut path to fully typed downstream occurrences.  Abstract quotient,
kernel and commuting-square tests are proved, but the actual lift and
all actual downstream maps remain unavailable.

The same batch contains a genuine arithmetic advance.  An exact
quotient identity places a fixed-two M\"obius periodic core in a
small-polylogarithmic almost-scale correlation theorem.  We retain
this as \(\Lone_{\rm actual\mbox{-}core}\) evidence.  It is not
positive \(\Ltwo\): the full physical scope and a fixed \(X\)-power
saving are absent.

We freeze the combined state in an MVP5 classifier.  Validation
precedes seven ordered route verdicts.  Architecture infeasibility
requires a theorem-backed complete route universe; rerouting requires
a fresh registry and typed crosswalk; a route stop requires exact
complete-cell scope; and the arithmetic frontier requires every
structural gate, an arithmetically independent H9 registry and a
strict physical-loss endpoint.  The current source-locked verdict is
\[
 \boxed{\operatorname{Verdict}_{143:151}=\textsc{not-testable}},
\]
with minimal missing set
\(\{\mathsf{H1.frontier\_occurrence\_lift}\}\).
This is forward structural and actual-core progress, but no
hard-packet, prime-pair or twin-prime theorem.
\end{abstract}

\noindent\textbf{Keywords:} fixed shift; occurrence lift; route
decision; actual M\"obius core; arithmetic frontier; endpoint
certificate.

\medskip
\noindent\fcolorbox{red!55!black}{red!3}{%
\begin{minipage}{0.96\linewidth}
\textbf{Claim boundary.}
\(\NT\) means that an active typed artifact is unavailable; it does
not mean the selected mathematical route is false.  The stopped cells
below are narrower architectures, not a complete impossibility
theorem.  A real almost-scale M\"obius core theorem is not the full
physical H3 theorem.  No positive fixed-\(X\)-power \(\Ltwo\) result,
strict endpoint, \(\Bhard=o(X)\), prime-pair lower bound, or
twin-prime theorem is claimed.
\end{minipage}}

\section{From MVP4 to MVP5}

MVP4 moved the first-failure pointer from a missing native archive to
a frontier-totalization node after an exact three-way cut
\citep{WangTPC141,WangTPC142}.  TPC-143--146 refine that node into an
executable four-map contract \citep{WangTPC143,WangTPC146}.  Its first
object is
\[
 L_X:\mathbb C^{\mathcal C_X^{\rm ns}}
 \longrightarrow \mathbb C^{\mathcal O_X},
 \qquad {\bf1}^{T}L_X={\bf1}^{T},
\]
where
\[
 \mathcal C_X^{\rm ns}
 =
 \mathcal C_X^{\rm ETO}\sqcup\mathcal C_X^{\rm FUM}.
\]
The actual frozen sample has no ETO row and 2,988 FUM rows, but the
required domain remains both classes.

The structural batch proves:
\begin{enumerate}[label=(\roman*),leftmargin=2.2em]
\item \(P_{h_0}^{\rm cut}=I\), with an explicit prohibition against
  calling it the downstream selector;
\item field-descent, \(Q_D/Q_Z\) kernel and \(G/P_{h_0}\)
  occurrence-level commutation criteria;
\item a zero-defect contract for the completed four-map route; and
\item scoped stops for current-schema-only and
  aggregate-cancellation implementations.
\end{enumerate}
Here ``current-schema-only'' means the maximal formal completion class
constrained by the archived fields and identities.  It is not an
impossibility theorem on the unknown actual-support carrier.
It does not construct \(L_X\).  The independent scalar route is not
closed either: it requires both a complete original-frontier
\(o(X)\) theorem and a theorem-backed disposition of every ETO path.
Thus the empty ETO class in the finite fixture and a frontier scalar
estimate alone cannot certify all-nonsoft totalization.

TPC-149 independently proves an actual fixed-two M\"obius periodic
core almost-scale estimate, and TPC-150 records the missing physical
return and an atomic-prefix firewall \citep{WangTPC149,WangTPC150}.
TPC-151 integrates the two branches without a promotion edge
\citep{WangTPC151}.

\section{A valid source-locked snapshot}

\begin{definition}[Valid MVP5 snapshot]
A snapshot is valid when:
\begin{enumerate}[label=(\alph*),leftmargin=2em]
\item its TPC-151 manifest and schema have the recorded canonical
  \texttt{CANONICAL\_UTF8\_LF\_V2} hashes;
\item hashes are explicitly integrity-only;
\item every import matches status, level, role, direction, scope,
  carrier, normalization and coverage;
\item the proof DAG is acyclic and its displayed order is canonical;
\item the selected active set is derived from the selected-root
  ancestor closure;
\item the route and stop tables are exact typed covers of the
  \emph{declared} tested cells;
\item H9 has no direct or indirect arithmetic-role or H2--H5 ancestor;
  and
\item all rational endpoint records and claim-boundary Booleans pass
  their direction checks.
\end{enumerate}
\end{definition}

Failure of one of these conditions returns \(\INVALID\) before any
mathematical route verdict.  In particular, source drift makes the
certificate stale; it does not refute a theorem.

\section{The minimal missing antichain}

Let \(\mathcal A\) be the active ancestor closure of the selected
root, including the root.  Define
\[
 \mathcal M=
 \{v\in\mathcal A:
   \operatorname{status}(v)=\NT
   \text{ or }v\text{ fails scope matching}\}.
\]
The partial-order frontier of missing nodes is
\[
 \min(\mathcal M)=
 \{v\in\mathcal M:
   \text{no ancestor of }v\text{ belongs to }\mathcal M\}.
\]
Only for deterministic reporting, the first representative is chosen
by canonical topological rank and node identifier.

\begin{proposition}[Frozen MVP5 missing frontier]
\label{prop:missing}
For the selected occurrence-lift/quotient-M\"obius route,
\[
 \boxed{
 \min(\mathcal M)=
 \{\mathsf{H1.frontier\_occurrence\_lift}\}.}
\]
\end{proposition}

\begin{proof}
The cut selector and abstract interface nodes are proved.  The
occurrence lift is \(\NT\).  Every actual \(Q_D,Q_Z,G\), downstream
\(P_{h_0}\), H1 carrier, H5--H9 and physical-return missing node is a
descendant of it in the selected DAG.  It is therefore the unique
minimal missing node.
\end{proof}

This algorithm prevents a later missing map from hiding its missing
parent and avoids treating incomparable prerequisites as if they had
a canonical mathematical order.

\section{Eight outcomes and their precedence}

The outer classifier has one validation outcome and seven
valid-snapshot verdicts.

\begin{definition}[Ordered valid-snapshot verdicts]
After validation, test:
\begin{enumerate}[label=(\roman*),leftmargin=2.2em]
\item \(\GO\): all active required gates are proved in compatible
  scopes and the full-synthesis endpoint has strict positive slack.
\item \(\AINF\): an independent theorem proves the declared route
  universe complete and exact complete-cell negative certificates
  stop every route.
\item \(\REROUTE\): the selected route is stopped and a distinct open
  typed alternative, fresh registry and theorem-backed crosswalk are
  selected.
\item \(\STOPR\): the selected route has a complete exact-cell stop
  certificate and no typed alternative is selected.
\item \(\NT\): an active required artifact is unavailable or
  scope-incompatible.
\item \(\AF\): every structural gate is proved, H9 is arithmetically
  independent, the physical-loss endpoint passes strictly, and every
  unresolved active node is an open scope-matched positive-\(\Ltwo\)
  arithmetic target.
\item \(\OPEN\): every remaining valid case.
\end{enumerate}
\end{definition}

\begin{theorem}[Total deterministic classifier]
Validation followed by this precedence returns exactly one of
\[
\begin{gathered}
 \INVALID,\quad \GO,\quad \AINF,\quad \REROUTE,\\
 \STOPR,\quad \NT,\quad \AF,\quad \OPEN.
\end{gathered}
\]
\end{theorem}

\begin{proof}
Validation either fails or provides a well-typed finite snapshot.
On a valid snapshot the ordered list ends in a catch-all, so its first
true predicate is unique.  The precedence deliberately separates a
complete architectural stop, a selected-cell stop, a missing object
and the arithmetic frontier.
\end{proof}

\section{Strengthened stop and reroute rules}

A list of tested routes is not automatically the universe of every
possible argument.  Therefore
\[
 \text{all declared cells stopped}
 \;\not\Longrightarrow\;
 \AINF
\]
unless an independent source proves the list complete in the actual
physical scope.

Likewise a negative result is stored with its exact scope, carrier,
normalization, registry and coverage.  The schema-only quotient
obstruction does not stop an occurrence-augmented quotient route; an
aggregate-commutator counterexample does not stop a row-separated
lift; and an atomic-prefix obstruction does not stop every possible
all-prefix theorem.

Finally, a reroute is not a change of a string identifier.  The
alternative must be open, use a distinct registry and possess a
theorem-backed crosswalk.  Otherwise the snapshot is invalid rather
than \(\REROUTE\).

\section{Three evidence labels}

The batch requires three progress labels:
\[
\begin{array}{ll}
\Lone_{\rm structural}:&
 \text{physical interfaces, typed covers and completion contracts},\\
\Lone_{\rm actual\mbox{-}core}:&
 \text{a genuine arithmetic theorem in a strict core scope},\\
\Ltwo_{\rm actual\mbox{-}positive}:&
 \text{a fixed-\(X\)-power saving on the literal physical carrier}.
\end{array}
\]

The TPC-149 theorem belongs to the middle line.  It evaluates actual
M\"obius factors after an exact quotient lift and is stronger than a
synthetic model.  It remains below the third line because its scope is
almost-scale and periodic-core, it lacks the complete physical return,
and its fixed-\(X\)-power exponent is zero.

\begin{proposition}[No pseudo-frontier promotion]
A conditional return formula, an \(\Lone_{\rm actual\mbox{-}core}\)
theorem, a frozen-family estimate or a logarithmic saving cannot be
the only unresolved evidence admitted by the \(\AF\) predicate.
\end{proposition}

\begin{proof}
The predicate requires unresolved records to have status
\(\OPEN\), evidence type \(\Ltwo_{\rm target\mbox{-}positive}\), exact
actual scope and nonstructural role.  Each listed substitute fails at
least one required field.
\end{proof}

\section{The split \texorpdfstring{\(1/400\)}{1/400} endpoint}

MVP5 separates:
\[
\begin{array}{rcl}
\text{arithmetic target}&:&
 \sigma_{\rm raw}\ge1/400,\\
\text{physical-loss target}&:&
 \Lambda_{\rm phys}<1/400,\\
\text{full synthesis}&:&
 \sigma_{\rm raw}-\Lambda_{\rm phys}>0.
\end{array}
\]
The \(\AF\) verdict requires the second line but deliberately leaves
the first as an open arithmetic target.  The \(\GO\) verdict requires
all three lines in compatible scope.

For the physical ledger, an upper certificate
\(U<1/400\) is a strict pass.  A complete lower or exact dual
certificate \(L\ge1/400\), including equality, stops the corresponding
route cell.  A conservative upper certificate \(U\ge1/400\) without a
lower certificate is neither pass nor stop.  Unknown tokens are not
zero, energy exponents are halved once, joint replacements cover
disjoint primitive occurrences, selector cost is charged once and
determinant reserve is not physical slack.

The current H9 registry is incomplete, so the physical and full
ledgers are both incomplete.  The log-power actual-core saving is
recorded separately at fixed-\(X\)-power exponent zero.

\section{H9 arithmetic independence}

The H9 occurrence registry is a structural/physical-synthesis
certificate.  Its allowed ancestors are H1, H6, H7, H8 and independent
physical cost sources.  The audit first rejects any ancestor whose
role is arithmetic shadow, arithmetic target or arithmetic negative,
then independently rejects H2--H5 gate ancestors.

This two-key check matters because a node could otherwise be renamed
outside an H3 namespace while still importing the target
correlation.  In particular, the actual-core theorem cannot be used
to fill an unknown physical occurrence cost, and H5 determinant
reserve cannot be registered as endpoint slack.

\section{Frozen gate projection}

\begin{table}[htbp]
\centering
\small
\caption{MVP5 gate projection.}
\label{tab:gates}
\begin{tabularx}{\textwidth}{@{}P{0.07\textwidth}P{0.19\textwidth}
P{0.21\textwidth}Y@{}}
\toprule
Gate & Status & Evidence target & First unresolved content\\
\midrule
H1 & \(\NT\) & \(\Lone_{\rm structural}\) &
Conservative row-separated occurrence lift on ETO and FUM.\\
H2 & \(\OPEN\) & positive \(\Ltwo\) &
Literal signed resonance replacement.\\
H3 & \(\OPEN\) & positive \(\Ltwo\) &
Physical weight, generic phase, all-prefix selector, four-sign return,
and a fixed-\(X\)-power saving.\\
H4 & \(\OPEN\) & positive \(\Ltwo\) &
Complete original-scale high/ultra/boundary return.\\
H5 & \(\NT\) & positive \(\Ltwo\) &
Actual determinant/zero pair on completed occurrence fibers.\\
H6 & \(\NT\) & \(\Lone_{\rm structural}\) &
Complete physical occurrence cover.\\
H7 & \(\NT\) & \(\Lone_{\rm structural}\) &
Actual downstream fixed-\(h_0\) totality.\\
H8 & \(\NT\) & \(\Lone_{\rm structural}\) &
Exact hard-packet reconnection.\\
H9 & \(\NT\) & \(\Lone_{\rm structural}\) &
Complete independent registry and physical-loss upper certificate.\\
\bottomrule
\end{tabularx}
\end{table}

The selected route is
\nolinkurl{occurrence_lift_quotient_mobius_return}.  It has no stop
certificate.  Six narrower cells have typed negative records, but the
route-universe completeness status is \(\OPEN\).

\section{MVP5 decision}

\begin{theorem}[Source-locked MVP5 verdict]
\label{thm:decision}
For the canonical TPC-143--151 snapshot dated 27 July 2026,
\[
 \boxed{
 \operatorname{Verdict}_{143:151}=\NT,
 \qquad
 \min(\mathcal M)=
 \{\mathsf{H1.frontier\_occurrence\_lift}\}.}
\]
The verdict is neither \(\GO\), \(\AF\), \(\AINF\), nor a route stop.
\end{theorem}

\begin{proof}
The canonical source and type checks pass, so the snapshot is not
\(\INVALID\).  The selected route is open, and no theorem proves the
declared route list complete, excluding all stop verdicts.  By
\cref{prop:missing}, an active structural artifact is unavailable, so
\(\NT\) occurs before \(\AF\) and \(\OPEN\).  Independently, H6--H9
and the physical endpoint are incomplete, while TPC-149 supplies
\(\Lone_{\rm actual\mbox{-}core}\), not positive physical \(\Ltwo\).
\end{proof}

\begin{corollary}[What moved]
The first-failure pointer moved from generic frontier totalization to
one explicit sparse matrix with an exact domain, conservation law and
lineage schema.  The arithmetic evidence simultaneously moved from a
more remote affine shadow to an actual M\"obius periodic core.  These
are forward moves, but neither closes a fixed-\(h_0\) physical
positive-\(\Ltwo\) gate.
\end{corollary}

\section{Regression suite and next forced step}

The standard-library audit reaches all seven valid verdicts on
synthetic typed snapshots and rejects:
\begin{enumerate}[leftmargin=2em]
\item a malformed route before verdict selection;
\item current-schema nonidentifiability promoted to actual-carrier
  impossibility;
\item architecture infeasibility without a route-universe theorem;
\item rerouting without a fresh registry and typed crosswalk;
\item conditional or actual-core pseudo-frontiers;
\item a stale or descendant-pre-empted first-missing pointer;
\item a frontier-totalization node whose artifact drops the ETO
  clause from the scalar route;
\item the positive-\(\Ltwo\) target label reported as an achieved
  fixed-power result;
\item an actual-core log saving assigned a positive \(X\)-power;
\item direct or indirect arithmetic ancestors of H9; and
\item a physical or full endpoint pass with missing registry data.
\end{enumerate}

On the selected occurrence-augmented map route, the next forced
structural object is \(L_X\), not another refinement of an aggregated
scalar.  A successful construction must cover both ETO and FUM,
conserve every source column, record every target occurrence
separately and preserve native, shift, normalization and support
provenance.  It would then make the four-map defect vector testable.
The distinct scalar-plus-ETO route remains open under its own complete
two-clause contract.

On the arithmetic side, the next separate task is to return the actual
periodic core through the literal physical weights, phases and
four-sign packet with either a deterministic all-prefix theorem or a
valid non-atomic selector.  Only after the structural endpoint passes
independently can the snapshot legitimately return \(\AF\).

\paragraph{Established.}
The valid-snapshot calculus, strengthened route-stop rules,
minimal-missing antichain, evidence classification, H9 firewall and
split endpoint semantics are \(\Lzero/\Lone\) results.

\paragraph{Not established.}
No occurrence lift, complete four-map physical carrier, positive
fixed-\(X\)-power \(\Ltwo\) saving, endpoint pass,
\(\Bhard=o(X)\), prime-pair lower bound or twin-prime theorem is
proved.

\section{Conclusion}

MVP5 keeps both gains visible without combining them into a claim that
neither source proves.  The structural route now has one concrete
first matrix, and the arithmetic route has one real actual-core
theorem.  The bridge between them is still missing.  The honest
source-locked verdict is therefore \(\NT\), with a sharper and more
actionable pointer than MVP4.

\bibliographystyle{plainnat}
\bibliography{references}

\end{document}
